\section{Introduction}
\label{sec:introduction}

% Modern language-model RL is implemented by two software stacks. An inference engine such as vLLM or SGLang generates rollouts, while a training engine such as Megatron or FSDP recomputes log probabilities and backpropagates gradients. Different kernels, parallelism layouts, reduction orders, precisions, and decode paths make these stacks numerically distinct even when they load the same parameters \citep{yao2025tis,he2025nondeterminism,qi2025fp16}. The discrepancy is especially consequential in mixture-of-experts models. A small numerical perturbation can change a discrete top-$k$ routing decision, so the training path may traverse experts that did not generate the sampled token \citep{ling2025ring,ma2025r3}. The mismatch can then affect both the token probability and the conditional mean of the gradient payload.

% Practice has produced two different repairs. Truncated importance sampling (TIS) caps the engine ratio \citep{ionides2008truncated,yao2025tis}; IcePop zeroes every ratio outside a fixed interval \citep{ling2025ring}. Both require choices outside the underlying RL objective. This leaves three questions unanswered: why should a large ratio be capped rather than removed, when should it instead be removed, and what uncertainty assumptions make either operation uniformly valid?

% \paragraph{Our view.}
% An uncertainty class is not merely a regularizer around a fixed algorithm. Together with a decision criterion, it selects the token-level operator and determines the boundary of its guarantee. A bounded-ratio class treats every payload as valid but potentially influential, leading to a projection onto an influence budget. A localized corrupted channel allows a small fraction of payload means to be unknown without imposing a zero-mean condition, leading to a clean-fidelity versus corrupted-exposure tradeoff. The cap and the gate are therefore solutions to different statistical problems rather than interchangeable heuristics.

% Our contributions are threefold.
% \begin{itemize}
%     \item We formalize geometry--operator correspondence through two requirements: an operator must minimize a bias constant uniformly over a stated observation-law class, and a matching hard instance must characterize the boundary beyond which its guarantee fails. This follows the closed-form-plus-lower-bound proof organization of \citet{shi2023distributionally}, while changing the statistical object from a robust Bellman functional to a ratio-weighted gradient estimator.
%     \item We prove a complete density-band correspondence. A cap covering the band is exactly unbiased; under a fixed influence budget it is the unique pointwise projection for every admissible ratio law; and an essential escape component places the observation law outside every finite band and creates a bias floor of at least $\eps B e^{\sigma}$ for every band-faithful cap.
%     \item We prove a complete localized-corruption correspondence. Without localization, the cap-versus-gate choice is not identifiable from the class. With localization, the uniform-risk optimal rule is the identity outside the escape set and $\min\{k,m^{\star}\}$ inside it. Hard zeroing is optimal if and only if positive clean escape mass is at most $\eps/(1-\eps)$, and an attaining-member construction matches the uniform constant for every measurable operator.
% \end{itemize}

\paragraph{Notation.}
We use plain letters such as $x$ to denote scalars, lowercase bold letters such as
$\mathbf{x}$ to denote vectors, and calligraphic letters such as $\mathcal{F}$ to
denote sets and function classes. For a vector $\mathbf{x}$, $\|\mathbf{x}\|_2$
denotes its $\ell_2$-norm. For a scalar $x \in (0,1)$, $\logit x := \log\frac{x}{1-x}$.